% Source: https://artofproblemsolving.com/community/c3h2271287_factoring_a_symmetric_polynomial
% Stable identity: aops-topic-2271287-q01
\textbf{Step 1: the factor $z+1$.} The coefficients come in equal neighbouring pairs, so
\[p(z)=(z+1)\left(4z^{10}-21z^8+17z^6+17z^4-21z^2+4\right).\]

\textbf{Step 2: put $w=z^2$.} The second factor is
\[r(w)=4w^5-21w^4+17w^3+17w^2-21w+4,\]
whose coefficients read the same backwards (a palindromic polynomial). A palindromic polynomial of odd degree has the root $w=-1$ (here $-4-21-17+17+21+4=0$), and division gives
\[r(w)=(w+1)\left(4w^4-25w^3+42w^2-25w+4\right).\]

\textbf{Step 3: the palindromic quartic.} Divide by $w^2$ and put $u=w+\dfrac{1}{w}$, so that $w^2+\dfrac{1}{w^2}=u^2-2$:
\[4\left(w^2+\dfrac{1}{w^2}\right)-25\left(w+\dfrac{1}{w}\right)+42=4u^2-25u+34=(u-2)(4u-17).\]
Multiplying back by $w^2$: $u-2$ gives $w^2-2w+1=(w-1)^2$ and $4u-17$ gives $4w^2-17w+4=(4w-1)(w-4)$. Hence
\[r(w)=(w+1)(w-1)^2(4w-1)(w-4).\]

\textbf{Step 4: back to $z$.} With $w=z^2$: $w+1=z^2+1$, $(w-1)^2=(z-1)^2(z+1)^2$, $4w-1=(2z-1)(2z+1)$ and $w-4=(z-2)(z+2)$. Altogether
\[\boxed{(z+1)^3(z-1)^2(z-2)(z+2)(2z-1)(2z+1)\left(z^2+1\right)}.\]
Each factor is irreducible over the integers ($z^2+1$ has no real roots), and the degrees add up to $3+2+1+1+1+1+2=11$.
